\subsection{COMMUTATORS IN A TOPOLOGICAL GROUP}

Let \(G\) be a topological group. We define the groups \(\overline{D^0}G, \overline{D^1}G, \overline{D^2}G, \ldots\) and \(\overline{C^1}G, \overline{C^2}G, \overline{C^3}G, \ldots\) by the formulae

\[
\begin{array}{l l} \overline {{\mathrm{D} ^ {0}}} \mathrm{G} = \mathrm{G}, & \overline {{\mathrm{D} ^ {i + 1}}} \mathrm{G} = (\overline {{\overline {{\mathrm{D} ^ {i}}} \mathrm{G}}}, \overline {{\mathrm{D} ^ {i}}} \mathrm{G}) \\ \overline {{\mathrm{C} ^ {1}}} \mathrm{G} = \mathrm{G}, & \overline {{\mathrm{C} ^ {i + 1}}} \mathrm{G} = (\overline {{\mathrm{G}}}, \overline {{\mathrm{C} ^ {i}}} \mathrm{G}). \end{array}
\]

PROPOSITION 1. Let G be a topological group and A and B subgroups of G. Then \((\overline{\mathrm{A}},\overline{\mathrm{B}}) = \overline{(\overline{\mathrm{A}},\overline{\mathrm{B}})},\overline{\mathrm{D}^{i}}\overline{\mathrm{A}} = \overline{\mathrm{D}^{i}\mathrm{A}},\overline{\mathrm{C}^{i}}\overline{\mathrm{A}} = \overline{\mathrm{C}^{i}\mathrm{A}}\) .

Let \(\phi\) be the continuous mapping \((x,y)\mapsto x^{-1}y^{-1}xy\) of \(\mathbf{G}\times \mathbf{G}\) into \(\mathbf{G}\) . Then \(\phi (\mathbf{A}\times \mathbf{B})\subset (\mathbf{A},\mathbf{B})\) , hence \(\phi (\overline{\mathbf{A}}\times \overline{\mathbf{B}})\subset (\overline{\mathbf{A}},\overline{\mathbf{B}})\) and hence

\[
\overline{(\overline{\mathrm{A}},\overline{\mathrm{B}})} \subset (\overline{\mathrm{A}},\overline{\mathrm{B}});
\]

the opposite inclusion is obvious and therefore \((\overline{\mathrm{A}},\overline{\mathrm{B}}) = \overline{(\overline{\mathrm{A}},\overline{\mathrm{B}})}\) . Clearly \(\overline{\mathrm{D}^{0}}\overline{\mathrm{A}} = \overline{\mathrm{D}^{0}\mathrm{A}}\) ; assuming the equality \(\overline{\mathrm{D}^{i}}\overline{\mathrm{A}} = \overline{\mathrm{D}^{i}\mathrm{A}}\) , it follows that

\[
\overline{\mathrm{D}^{i+1}}\overline{\mathrm{A}} = (\overline{\overline{\mathrm{D}^{i}}\overline{\mathrm{A}}},\overline{\mathrm{D}^{i}}\overline{\mathrm{A}}) = (\overline{\mathrm{D}^{i}}\overline{\mathrm{A}},\overline{\mathrm{D}^{i}}\overline{\mathrm{A}}) = \overline{(\mathrm{D}^{i}\mathrm{A},\mathrm{D}^{i}\mathrm{A})} = \overline{\mathrm{D}^{i+1}\mathrm{A}}
\]

and hence \(\overline{\mathrm{D}^{i}}\overline{\mathrm{A}} = \overline{\mathrm{D}^{i}\mathrm{A}}\) for all \(i\) . The proof of the formula \(\overline{\mathrm{C}^{i}}\overline{\mathrm{A}} = \overline{\mathrm{C}^{i}\mathrm{A}}\) is analogous.

COROLLARY 1. If G is Hausdorff, the following conditions are equivalent:

\begin{enumerate}
\def\labelenumi{(\roman{enumi})}
\item
  G is solvable (resp. nilpotent);
\item
  \(\overline{\mathrm{D}^{i}}\mathrm{G} = \{e\}\) (resp. \(\overline{\mathrm{C}^{i}}\mathrm{G} = \{e\}\) ) for sufficiently large \(i\) .
\end{enumerate}

\(\mathrm{D}^{i}\mathrm{G} \subset \overline{\mathrm{D}^{i}}\mathrm{G}, \mathrm{C}^{i}\mathrm{G} \subset \overline{\mathrm{C}^{i}}\mathrm{G}\) and hence (ii) \(\Rightarrow\) (i). \(\{e\} = \overline{\{e\}}\) and hence (i) \(\Rightarrow\) (ii) by Proposition 1.

COROLLARY 2. Let G be a Hausdorff topological group and A a subgroup of G. For A to be solvable (resp. nilpotent, commutative), it is necessary and sufficient that \(\overline{\mathbf{A}}\) be so.

This follows immediately from Proposition 1.

PROPOSITION 2. Let \(\mathbf{G}\) be a topological group and \(\mathbf{A}\) and \(\mathbf{B}\) subgroups of \(\mathbf{G}\) . If \(\mathbf{A}\) is connected, (A, B) is connected.

For fixed y in B, the set \(M_{y}\) of \((x,y)\) with \(x \in A\) is connected (for the mapping \(x \mapsto (x,y)\) of A into G is continuous). \(e \in M_{y}\) and hence the union R of the \(M_{y}\) with \(y \in B\) is connected. But (A, B) is the subgroup of G generated by R, whence the proposition.

